\documentclass{article}
\usepackage{amsmath, amssymb, amsthm}
\usepackage{hyperref}
\usepackage{geometry}
\geometry{margin=1in}

\title{Erd\H{o}s Problem \#247}
\date{Last updated: 2025-08-31}
\author{Source: erdosproblems.com}

\begin{document}
\maketitle

\noindent\textbf{Status:} OPEN \quad \textbf{No prize}

\noindent\textbf{Tags:} number theory, irrationality

\medskip
\noindent\textbf{Problem Statement:}

Let $1\leq a_1<a_2<\cdots$ be a sequence of integers such that\[\limsup \frac{a_n}{n}=\infty.\]Is\[\sum_{n=1}^\infty \frac{1}{2^{a_n}}\]transcendental?

\bigskip
\noindent\textbf{Remarks:}

Erd\H{o}s \cite{Er75c} proved the answer is yes under the stronger condition that $\limsup n_k/k^t=\infty$ for all $t\geq 1$.

Erd\H{o}s \cite{Er88c} says 'many of these problems seem hopeless at present, but perhaps one can prove that if $a_n>cn^2$ then $\sum_{n=1}^\infty \frac{1}{2^{a_n}}$ is not the root of any quadratic polynomial'.

This problem has been formalised in Lean as part of the Google DeepMind Formal Conjectures project.

\bigskip
\noindent\textbf{References:}

[Er75c] Erd\H{o}s, P., Some problems and results on the irrationality of the sum of infinite series. J. Math. Sci. (1975), 1-7 (1976).
 
 [Er88c] Erd\H{o}s, P., On the irrationality of certain series: problems and results. New advances in transcendence theory (Durham, 1986) (1988), 102-109.

\bigskip
\noindent\small{Source: \url{https://www.erdosproblems.com/247}}
\end{document}
